\documentclass[11pt]{article}
\usepackage[a4paper,margin=1in]{geometry}
\usepackage{fontspec}
\usepackage[boldfont=false]{xeCJK}
\setCJKmainfont{FandolSong-Regular.otf}
\usepackage{amsmath}
\usepackage{amssymb}
\usepackage{array}
\usepackage{booktabs}
\usepackage{graphicx}
\usepackage{adjustbox}
\usepackage{enumitem}
\setlist[itemize]{topsep=2pt,partopsep=0pt,parsep=0pt,itemsep=2pt,leftmargin=1.4em}
\setlist[enumerate]{topsep=2pt,partopsep=0pt,parsep=0pt,itemsep=2pt,leftmargin=1.6em}
\usepackage{parskip}
\usepackage{fvextra}
\fvset{breaklines=true,breakanywhere=true,fontsize=\small}
\setlength{\parindent}{0pt}

\begin{document}

\begin{center}
  {\LARGE\bfseries Understanding business activity}\\[0.35em]
  {\large Igcse Business Studies}
\end{center}
\vspace{0.6em}

\section*{What business activity is for}

A \textbf{business} 企业 is an organisation that makes \textbf{goods} 商品 or provides \textbf{services} 服务 for \textbf{customers} 顾客.

\begin{itemize}
\item Goods are physical things you can touch — bread, phones, cars.

\item Services are useful actions done for you — a haircut, a bus ride, a doctor's visit.

\end{itemize}

The purpose of business activity is to use \textbf{resources} 资源 to make goods and services, so that people's \textbf{needs} 需要 and \textbf{wants} 欲望 are met.

\begin{itemize}
\item Needs are things you must have to live: food, water, shelter, clothes.

\item Wants are things you would like but can live without: a holiday, a games console.

\end{itemize}

\subsection*{Factors of production}

To produce anything, a business brings together four \textbf{factors of production} 生产要素 — the resources needed to make goods and services.

\begin{center}
\adjustbox{max width=\linewidth}{%
\begin{tabular}{lll}
\toprule
\textbf{Factor} & \textbf{What it means} & \textbf{Example} \\
\midrule
\textbf{land} 土地 & natural resources & farmland, water, oil, minerals \\
\textbf{labour} 劳动力 & the work of people & workers, managers \\
\textbf{capital} 资本 & money and the man-made tools bought with it & machines, factories, vehicles \\
\textbf{enterprise} 企业家才能 & the skill to bring the other three together and take risks & the owner's idea and effort \\
\bottomrule
\end{tabular}}
\end{center}

\subsection*{The economic problem}

People's needs and wants are unlimited. But resources are \textbf{scarce} 稀缺 — there are not enough of them for everything. This is the \textbf{economic problem} 经济问题.

Because resources are scarce, you must always choose. You cannot have everything.

\subsection*{Opportunity cost}

When you choose one thing, you give up the next best thing. The \textbf{opportunity cost} 机会成本 is the next best choice you give up.

Example: a business has some money and spends it on a new machine. Now it cannot also spend that money on a bigger shop. The bigger shop is the opportunity cost.

\subsection*{Added value}

\textbf{Added value} 增值 is the difference between the price a business sells a product for and the cost of the \textbf{raw materials} 原材料 bought in to make it.

\[
\text{added value} = \text{selling price} - \text{cost of bought-in materials}
\]

A business can increase added value by:

\begin{itemize}
\item raising the price, often through a strong \textbf{brand} 品牌 or better design,

\item lowering the cost of materials, for example by finding cheaper suppliers,

\item improving the product so customers will pay more.

\end{itemize}

\subsection*{Why businesses succeed or fail}

A business is more likely to succeed when it has:

\begin{itemize}
\item good \textbf{management} 管理 — skilled leaders who plan well,

\item enough \textbf{finance} 资金 — money to pay bills while it grows,

\item a product customers want, and an edge over the \textbf{competition} 竞争 (the other businesses selling similar things).

\end{itemize}

A business may fail and close down when it runs out of cash, manages money badly, or cannot compete. Weak management and poor cash control are the most common causes.

\section*{Classifying businesses}

\subsection*{The three sectors}

We sort all business activity into three groups, by how far the work is from the final customer.

\begin{center}
\adjustbox{max width=\linewidth}{%
\begin{tabular}{lll}
\toprule
\textbf{Sector} & \textbf{What it does} & \textbf{Examples} \\
\midrule
\textbf{primary sector} 第一产业 & takes natural resources from the earth & farming, fishing, mining, oil drilling \\
\textbf{secondary sector} 第二产业 & makes goods from raw materials & building, \textbf{manufacturing} 制造业, food processing \\
\textbf{tertiary sector} 第三产业 & provides services & shops, banks, transport, tourism \\
\bottomrule
\end{tabular}}
\end{center}

\subsection*{How the sectors change}

The importance of each sector changes as a country grows richer.

\begin{itemize}
\item In a \textbf{developing economy} 发展中经济体 (a poorer country building up its industry), the primary sector is often large — many people farm.

\item As the country grows, the secondary sector grows too.

\item In a \textbf{developed economy} 发达经济体 (a rich, industrialised country), most people work in the tertiary sector.

\end{itemize}

\textbf{De-industrialisation} 去工业化 is when the secondary sector shrinks and the tertiary sector grows. This has happened in many developed economies as factories move abroad.

\subsection*{Private sector and public sector}

The \textbf{private sector} 私营部门 is made up of businesses owned by private people and groups. Most aim to make a \textbf{profit} 利润. The \textbf{public sector} 公共部门 is made up of organisations owned by the government, which aim to provide a service to the public.

\begin{center}
\adjustbox{max width=\linewidth}{%
\begin{tabular}{lll}
\toprule
\textbf{} & \textbf{Private sector} & \textbf{Public sector} \\
\midrule
Owner & private people and companies & the government \\
Main aim & usually profit & provide services to everyone \\
Examples & shops, factories, private firms & state schools, public hospitals, the army \\
\bottomrule
\end{tabular}}
\end{center}

\section*{Enterprise, business growth and size}

\subsection*{The entrepreneur}

An \textbf{entrepreneur} 企业家 is a person who sets up a business and takes the risk of running it. Successful entrepreneurs usually:

\begin{itemize}
\item are willing to take \textbf{risk} 风险 — to risk their own money on a new idea,

\item work hard and do not give up,

\item are \textbf{innovative} 创新 — full of new ideas,

\item are confident and good at making decisions,

\item are well organised, able to manage time, money and people.

\end{itemize}

\subsection*{The business plan}

A \textbf{business plan} 商业计划书 is a document that describes the business and its future. It usually contains the business idea and objectives, the product, the \textbf{market} 市场 it will sell to, who will run it, and a plan for the money.

A business plan has two main uses:

\begin{itemize}
\item it makes the owner think carefully and plan ahead,

\item it helps the business \textbf{raise finance} 筹集资金 — banks and investors want to see a plan before they lend or invest.

\end{itemize}

\subsection*{Government support for start-ups}

A \textbf{start-up} 初创企业 is a brand-new business. New businesses create jobs and help the economy grow, so governments often support them with:

\begin{itemize}
\item \textbf{grants} 补助金 — money given that does not have to be paid back,

\item low-interest \textbf{loans} 贷款 — money lent on easy terms,

\item training and advice for new owners,

\item lower taxes in the first years.

\end{itemize}

\subsection*{Measuring the size of a business}

You can measure how big a business is in several ways. Each one has a weakness.

\begin{center}
\adjustbox{max width=\linewidth}{%
\begin{tabular}{lll}
\toprule
\textbf{Measure} & \textbf{How it works} & \textbf{Limitation} \\
\midrule
number of \textbf{employees} 员工 & count the staff & a high-tech firm can have few staff but large \textbf{output} 产量 \\
\textbf{revenue} 收入 (sales) & the value of everything sold & a firm selling cheap goods may sell a lot yet be small \\
\textbf{capital employed} 所用资本 & the total money invested in the business & hard to compare between different industries \\
\bottomrule
\end{tabular}}
\end{center}

\subsection*{Staying small or growing}

Some businesses stay small on purpose: the owner wants to keep control, the market is small, or the product is specialised.

Businesses that want to grow bigger — to achieve \textbf{growth} 增长 — can do it in two ways.

\begin{itemize}
\item \textbf{Internal growth} 内部增长 (also called organic growth): the business grows by itself, by selling more, opening new branches, or adding new products. It is slower but lower-risk.

\item \textbf{External growth} 外部增长: the business joins with another business. This is faster. It happens through a \textbf{merger} 合并 (two businesses agree to join into one) or a \textbf{takeover} 收购 (one business buys control of another).

\end{itemize}

External growth is classified by the kind of business that is joined.

\begin{center}
\adjustbox{max width=\linewidth}{%
\begin{tabular}{lll}
\toprule
\textbf{Type} & \textbf{What joins} & \textbf{Example} \\
\midrule
\textbf{horizontal integration} 横向一体化 & two firms at the same stage of the same business & two car makers join \\
\textbf{vertical integration} 纵向一体化 & two firms at different stages of the same business & a car maker buys a tyre maker \\
\textbf{conglomerate integration} 混合一体化 & two firms in completely different businesses & a car maker buys a food company \\
\bottomrule
\end{tabular}}
\end{center}

\subsection*{Economies and diseconomies of scale}

As a business grows, its cost per unit can fall. These savings are called \textbf{economies of scale} 规模经济.

\begin{center}
\adjustbox{max width=\linewidth}{%
\begin{tabular}{ll}
\toprule
\textbf{Type} & \textbf{How the cost per unit falls} \\
\midrule
purchasing & buying large amounts of materials brings a lower price \\
financial & large firms can borrow money more cheaply \\
managerial & the firm can afford expert managers \\
technical & the firm can use large, efficient machines \\
\textbf{marketing} 营销 & the cost of \textbf{advertising} 广告 is spread over more goods \\
\bottomrule
\end{tabular}}
\end{center}

If a business grows too big, the cost per unit can start to rise again. These are \textbf{diseconomies of scale} 规模不经济. They happen because \textbf{communication} 沟通 in a huge firm gets worse, workers feel less important and work less hard, and managers find the firm hard to control.

\section*{Types of business organisation}

\subsection*{Liability}

\textbf{Liability} 责任 means who must pay the \textbf{debts} 债务 if a business cannot pay them itself.

\begin{itemize}
\item With \textbf{unlimited liability} 无限责任, the owner is personally responsible for all the debts. If the business fails, the owner may lose personal things, like a house or car, to pay them.

\item With \textbf{limited liability} 有限责任, the owner can only lose the money they put into the business. Their personal things are safe.

\end{itemize}

\subsection*{Incorporated and unincorporated}

\begin{itemize}
\item An \textbf{incorporated} 法人 business has a separate legal identity from its owners. The business itself can own things and be taken to court. Its owners get limited liability. Companies are incorporated.

\item An \textbf{unincorporated} 非法人 business is not separate from its owner in law. Sole traders and partnerships are unincorporated.

\end{itemize}

\subsection*{The four main types}

A \textbf{sole trader} 独资企业 is a business owned by one person.

\begin{center}
\adjustbox{max width=\linewidth}{%
\begin{tabular}{ll}
\toprule
\textbf{Advantages} & \textbf{Disadvantages} \\
\midrule
easy and cheap to set up & unlimited liability \\
the owner keeps all the profit & hard to raise finance \\
the owner makes all decisions & long hours, no one to share the work \\
private (accounts stay secret) & the business ends if the owner stops \\
\bottomrule
\end{tabular}}
\end{center}

A \textbf{partnership} 合伙企业 is owned by two or more people who share the work, the decisions and the profit. They usually sign a written agreement.

\begin{center}
\adjustbox{max width=\linewidth}{%
\begin{tabular}{ll}
\toprule
\textbf{Advantages} & \textbf{Disadvantages} \\
\midrule
more owners bring more capital & unlimited liability (in most cases) \\
work and ideas are shared & the profit is shared \\
easy to set up & partners can disagree \\
 & one partner's mistake affects everyone \\
\bottomrule
\end{tabular}}
\end{center}

A \textbf{private limited company} 私人有限公司 (often "Ltd") is owned by \textbf{shareholders} 股东, but its shares cannot be sold to the public. The owners are often family and friends.

\begin{center}
\adjustbox{max width=\linewidth}{%
\begin{tabular}{ll}
\toprule
\textbf{Advantages} & \textbf{Disadvantages} \\
\midrule
limited liability & more legal rules and paperwork \\
easier to raise capital than a sole trader & shares cannot be sold to the public \\
control stays in a small group & accounts must be shared with the government \\
\bottomrule
\end{tabular}}
\end{center}

A \textbf{public limited company} 公众有限公司 (often "plc") is a large company whose \textbf{shares} 股份 can be bought by the public on the \textbf{stock exchange} 证券交易所.

\begin{center}
\adjustbox{max width=\linewidth}{%
\begin{tabular}{ll}
\toprule
\textbf{Advantages} & \textbf{Disadvantages} \\
\midrule
can raise very large amounts of capital & expensive and complex to set up \\
limited liability & accounts are fully public \\
well known, so easy to borrow & anyone can buy shares, so there is a risk of takeover \\
 & the original owners may lose control \\
\bottomrule
\end{tabular}}
\end{center}

Shareholders may receive a share of the profit, called a \textbf{dividend} 股息.

\subsection*{Other forms of organisation}

\begin{itemize}
\item \textbf{Franchise} 特许经营: a person (the \textbf{franchisee} 加盟商) pays to use the name, products and methods of an existing successful business (the \textbf{franchisor} 授权商). The franchisee gets a known brand, training and lower risk, but must pay fees and a share of the profit and has less freedom. Many fast-food chains work this way.

\item \textbf{Joint venture} 合资企业: two businesses agree to share the cost, risk and profit of one project, while staying separate companies.

\item \textbf{Social enterprise} 社会企业: a business that trades to help society or the environment, not only to make profit. Most of its profit is put back into its social aims.

\end{itemize}

\subsection*{Public sector organisation}

A \textbf{public corporation} 国有企业 is a business owned and run by the government, such as a national rail or postal service. Its aims include giving an important service to everyone, not only making a profit.

\section*{Business objectives and stakeholder objectives}

\subsection*{Business objectives}

An \textbf{objective} 目标 is a goal that a business aims for. Common objectives are:

\begin{itemize}
\item \textbf{survival} 生存 — staying in business, especially when new or in hard times,

\item profit — earning more than it spends; the main aim of most private firms,

\item growth — getting bigger, to be safer and earn more,

\item \textbf{market share} 市场份额 — the share of total sales in a market that one business holds,

\item \textbf{customer service} 顾客服务 — keeping customers happy so they come back,

\item social objectives — helping society or the \textbf{environment} 环境.

\end{itemize}

\subsection*{Objectives of a social enterprise}

A social enterprise aims for three kinds of objective together:

\begin{itemize}
\item economic — earn enough money to survive and grow,

\item social — help people, for example by giving jobs to those who find work hard to get,

\item environmental — protect the planet, for example by cutting waste.

\end{itemize}

\subsection*{Stakeholders}

A \textbf{stakeholder} 利益相关者 is any person or group that is affected by a business or has an interest in it.

\begin{center}
\adjustbox{max width=\linewidth}{%
\begin{tabular}{ll}
\toprule
\textbf{Stakeholder} & \textbf{What they usually want} \\
\midrule
owners / shareholders & good profit and a rising dividend \\
employees & fair pay, safe work, secure jobs \\
customers & good quality at a fair price \\
\textbf{suppliers} 供应商 & to be paid on time, and regular orders \\
government & businesses to obey laws and pay \textbf{tax} 税收 \\
\textbf{local community} 当地社区 & jobs, but no pollution or noise \\
banks & to be repaid the money they lent \\
\bottomrule
\end{tabular}}
\end{center}

\subsection*{Why stakeholder objectives conflict}

Stakeholders often want different things, so their objectives can \textbf{conflict} 冲突. A business cannot fully please everyone at once.

\begin{itemize}
\item Owners want higher profit, but employees want higher pay. Higher pay lowers profit.

\item Owners want to cut costs, but the local community wants less \textbf{pollution} 污染. Cutting costs may create more pollution.

\item Customers want low prices, but suppliers want to be paid more.

\end{itemize}

Good managers try to balance these different objectives.


\end{document}
